\section{A spectral obstruction on a marked punctured double}
\label{sec:marked-double-exterior}

Let \(W\) be a smooth compact oriented contractible four-manifold
with connected boundary \(Y\). An orientation-preserving diffeomorphism
\(f:Y\to Y\) gives the oriented double \(W\cup_f(-W)\).
Van Kampen and Mayer--Vietoris give a smooth homotopy four-sphere.
Insert \([0,L]\times Y\) using the identity and \(f\) on its two
faces. Collared reparametrization preserves the given smooth type
and is fixed off the seam collars. Mark an interior coordinate ball
on each side, disjoint from these collars. Smooth metrics on the
punctured halves can have arbitrary prescribed product seam
germs and exact unit round product germs at their marked
punctures. Use the pulled-back seam metrics in the corresponding
face markings and interpolate to interior metrics in the complements
of the prescribed germs. The round ends admit conformally flat
caps in their original ball parametrizations.

On a closed four-manifold write \(\Lambda\) for the maximum of
the largest signed eigenvalues of the chiral Weyl blocks. The
ordinary and modified quadratic forms have integrand
\(6|dv|^2+Vv^2\), with \(V=R\) and \(V=R-6\Lambda\),
respectively. On the retained exterior the contribution from a
unit round end is
\(6\langle\tau,\sqrt{1-\Delta_{S^3}}\tau\rangle\) for its trace
\(\tau\), by minimizing on the added conformally flat half
cylinder. This follows by decomposing into eigenfunctions of
\(-\Delta_{S^3}\); each mode minimizes the one-dimensional
quadratic action with prescribed value and decay at infinity.
For the identity double with reflected metrics and corresponding
marked punctures, use all \(H^1\) tests on one punctured half
and leave the trace on the doubled seam free. Positivity of
the corresponding one-half form, including its round-end term,
is equivalent to positivity of the full exterior form.
Restriction gives the sufficient implication by addition,
and even extension gives necessity since it doubles the value.
Orientation reversal exchanges chiralities and preserves
\(\Lambda\). It also makes the full positive chiral energy
the sum of both half energies, with no Weyl energy in the caps.

\begin{proposition}[Conformally product seam obstruction]
\label{prop:marked-double-neck}
Suppose the retained metric on the whole strip \([0,L]\times Y\)
is conformal to \(dt^2+\gamma\), where \(\gamma\) is a smooth
metric on the closed three-manifold \(Y\). Let \(r_{\min}\) be
the least Ricci eigenvalue of \(\gamma\) and use the transverse
first eigenvalues
\[
 \ell_{\rm mod}=\lambda_1(-6\Delta_\gamma+3r_{\min}),\qquad
 \ell_{\rm ord}=\lambda_1(-6\Delta_\gamma+R_\gamma).
\]
If either \(\ell<0\) and \(L\sqrt{-\ell}\ge\pi\sqrt6\), the
corresponding exterior form fails strict positivity, regardless
of the half metrics, puncture caps or twist. In particular, the
ordinary test excludes positive Yamabe sign on every closed
completion that retains this strip. The strip contributes the
exact energy \(\frac L4\int_Y|\operatorname{Ric}_{0,\gamma}|^2
\,dV_\gamma\) to either chirality.
\end{proposition}

\begin{proof}
In an orthonormal product frame whose last three vectors
diagonalize \(\operatorname{Ric}_\gamma\), the only sectional
curvature in the opposite pair \(0i,jk\) is
\(K_{jk}=R_\gamma/2-r_i\). Project onto either chiral basis and
subtract the scalar contribution \(R_\gamma/12\).
The Weyl eigenvalues are \(R_\gamma/6-r_i/2\), giving the
modified potential and energy identity. The lowest Dirichlet
value of the product strip is \(\ell+6\pi^2/L^2\), by the
Poincaré inequality in \(t\) and the transverse Rayleigh
inequality. It is attained by \(\sin(\pi t/L)\phi(y)\) for a
first transverse eigenfunction \(\phi\). Its zero extension
is a nonzero \(H^1\) test on the retained exterior, with zero
trace on every marked round end. The scalar quadratic form
transfers under multiplication of the test by the inverse
conformal factor, since the test has zero trace on both strip
faces. The added potential transfers because \(\Lambda\) has
conformal weight \(-2\). The Weyl energy of the strip is
conformally invariant in four dimensions. A warped collar
\(ds^2+a(s)^2\gamma\) is included with \(L=\int ds/a(s)\).
\end{proof}

The same product calculation gives a shorter bound for fixed
half metrics. Normalize \(\phi\) in \(L^2(\gamma)\), take one
test on each half with that trace in the common cylinder
markings and zero trace at the marked punctures, and let \(H\)
be the sum of their form values off the inserted strip. Put
\(k=\sqrt{-\ell/6}\). When \(0<L<\pi/k\), the intervening test
\[
 \frac{\cos(k(t-L/2))}{\cos(kL/2)}\phi(y)
\]
has both required traces. Integration by parts gives strip
value \(-12k\tan(kL/2)\), so \(H\le12k\tan(kL/2)\) suffices
for failure. The bound includes the cost of realizing the
two traces in the fixed halves. In the second copy the
actual trace in its own marking is \(\phi\circ f^{-1}\).

\begin{example}[The Mazur seam]\label{ex:mazur-spectral-seam}
For \(n=1\), the lower half of
\href{https://arxiv.org/pdf/2403.17356v2#page=35}{Isoshima's
Figure 37} is the Mazur manifold \(W^-(0,3)\), with the cut
after its first two-handle. The identity double is standard
by Isoshima's Theorem 3.3. The punctures can be marked in
the two corresponding zero-handles. The actual middle
cut is \(Y=\Sigma(2,5,7)\), using
\href{https://math.berkeley.edu/~kirby/papers/Akbulut%20and%20Kirby%20-%20Mazur%20manifolds%20-%20MR0544597.pdf}{Akbulut--Kirby,
Proposition 1 and Theorem 2(1)}.

For every \(\gamma\) on this \(Y\), both transverse
eigenvalues in Proposition~\ref{prop:marked-double-neck}
are negative. Indeed, \(1/2+1/5+1/7<1\), so
\href{https://webhomes.maths.ed.ac.uk/~v1ranick/papers/milnbries.pdf}{Milnor,
Theorem 2.7 and Corollary 6.6}, gives a finite oriented
circle-bundle cover \(Z\) of \(Y\) over a closed orientable
surface of genus at least two. The preimage of an essential
simple closed curve is a torus with injective fundamental
group, by the homotopy sequence of the fibration. The
projection onto the surface group also gives \(b_1(Z)\ge4\).
Any nonnegative scalar curvature metric on \(Y\) would lift
to \(Z\) and force that metric to be flat by
\href{https://archive.ymsc.tsinghua.edu.cn/pacm_download/59/11113-Shing-Tung_Yau_216.pdf}{Schoen--Yau,
Theorem 3, p.~2567}. Parallel harmonic one-forms would give
\(b_1(Z)\le3\), a contradiction. Thus the lowest eigenvalue
of \(-8\Delta_\gamma+R_\gamma\) is negative, since a
nonnegative value would give a nonnegative scalar curvature
metric by conformal change with the positive first
eigenfunction \cite[eq.~(1.1)]{LeeParker1987Yamabe}. Both
transverse forms are bounded above on that test by this
three-dimensional conformal form.

For an explicit collar take the metric descending from the
left-invariant metric on
\(\widetilde{\mathrm{SL}}(2,\mathbb R)\) with cyclic
orthonormal bracket constants \((a,a,-a)\), \(a>0\).
Use Milnor's Brieskorn quotient model (Corollary 6.5)
and his
\href{https://core.ac.uk/download/pdf/82428733.pdf}{Ricci
formula (Theorem 4.3)}. The eigenvalues of Ricci are
\((-3a^2/2,-3a^2/2,a^2/2)\). The modified and ordinary
transverse potentials are constant, with the corresponding
obstructions at \(aL\ge 2\pi/\sqrt3\) and
\(aL\ge\pi\sqrt{12/5}\). Either chiral charge is
\((2/3)a^4L\operatorname{Vol}_{\gamma}(Y)\).
This supplies a restriction on the conformally product
construction at the specified marked cut. Shorter collars
require an estimate on the half metrics or further geometric
data.
\end{example}
